\section{Numerical methods}
\label{app:methods}

\paragraph*{Gaussian mixtures.} Mixtures with up to $K = 32$ components were searched for small $g$, small $P$ and large values of the ratios of the line route, with the focus constraint $\Upsilon \leq 1/2$ imposed exactly. Each final state was re-evaluated on two grids of different resolution, and only states satisfying the constraint on both were kept. The smallest gap found, $0.884\,\sigma$ at $\sigma = 10^{-2}$, belongs to an ``hourglass'': a narrow needle at the top with weaker components spread along $w$. It was confirmed by three forms of the gap formula and by an independent quadrature written separately from the search code (relative difference below $5\times 10^{-4}$).

\paragraph*{Closed forms for the strips.} Each term of Eq.~\eqref{eq:strip} is a Gaussian in $w$ times, in $v$, either a Gaussian or the real part of $N(v;0,L_v^2)\,e^{ibv} = e^{-b^2L_v^2/2}N(v;\,ibL_v^2,\,L_v^2)$ with $b = k_0/\sigma$, a Gaussian with a complex mean. Blurs and posteriors of such terms are Gaussian with explicit means and variances, so that $\varrho$, $s$, $H$ and the posterior moments in Eq.~\eqref{eq:Ncov} are closed-form sums. The covariance in Eq.~\eqref{eq:Ncov} is evaluated with Stein's identity within each term and an explicit between-term part, in which all differences of nearly equal quantities are formed analytically. The fast phase $k_0 v/\sigma$ is integrated as an independent variable over one period at fixed slow position, with error of order (period$/L_v)^2$.

\paragraph*{Pitfalls.} Three effects produced wrong values in the course of the study and are guarded against in the final codes.
\begin{enumerate}
\item \emph{Rounding.} In the limit $\sigma^2/L_v^2 \to 0$, the matrix $I - \sigma^2 T^{-1}$ rounds to the identity in double precision, and a derivative computed from it carried a rounding error of $2\times 10^{-17}$ against a true value of $5\times 10^{-21}$. An earlier version of the study reported a floor of the strip gap at $0.234\,\sigma$ because of it. The derivative is now computed as $-\sigma^2 e_z\cdot T^{-1}e_p$, and a pointwise evaluation in 320-digit arithmetic agrees to $3\times 10^{-14}$ at $\sigma = 10^{-80}$.
\item \emph{Phase precision.} Positions written as an absolute coordinate plus a fraction of a ripple period lose the phase once $L_v/\sigma \gtrsim 10^{12}$. In that regime only phase-based evaluators are used.
\item \emph{Envelope aliasing.} Averaging one period at coarse slow positions aliases an oscillation of relative size period$/L_v$; at $\sigma = 10^{-6}$ this was a $3\times 10^{-3}$ error, found because a full two-dimensional grid and a phase-based evaluator agreed with each other but not with the third method.
\end{enumerate}
A lesson for the cross-checks: in the stage that produced the artefact of item 1, two evaluations with different integration schemes agreed because they shared the routine that computed the pointwise fields. Agreement between evaluations is evidence only if they share no code for the states, the pointwise fields and the integration. All values reported here satisfy this.

\paragraph*{Independent evaluation.} The strip values at $\sigma = 10^{-4}$, $10^{-6}$, $10^{-12}$, $10^{-24}$ and $10^{-48}$ were reproduced by an evaluator written separately from the closed forms above, with no code in common with the study. Its formulas for $\varrho$, $s$, $D$ and $\partial_p D$ were checked pointwise against direct two-dimensional quadrature of the posterior (relative differences below $10^{-6}$). It reproduced $g$ to ten digits or better, and an optimisation with it from unrelated starting points found $g_{\min} = 0.33487$ at $\sigma = 10^{-12}$ and $0.17938$ at $\sigma = 10^{-24}$, with the same optimal parameters.

\paragraph*{Fields of the heat-flow route.} The quantities of Sec.~\ref{sec:removing} ($J_{ww}$, $R$, $t^2R'$, the eigenvalues of $H$) were computed on two-dimensional grids from closed-form fields of mixtures and lattices. For the strips the curvature does not depend on the slow coordinate, so the largest eigenvalue of $-tH$ was taken over the full plane of the fast phase (720 points) and $y_w$ (about 5000 points) at 25 blurs; the maximum is unchanged when the region is restricted to densities above $10^{-8}$ or $10^{-12}$ of the maximum, and when the grid is refined. An earlier evaluation that took the maximum over quadrature nodes only gave $1.82$ instead of $1.87$. Separately from the study's code, the values of $-tH$ for the strip at $\sigma = 10^{-2}$ were reproduced on a two-dimensional Fourier grid of the state.

\paragraph*{Verification suite.} The repository (Data and Code Availability) contains the results, the figure scripts and a verification suite of 83 checks that recompute every number and inequality of this article from the stored results, with light recomputations: the hourglass minimum in three forms, a 320-digit pointwise check, the strip optima by two code-disjoint evaluators, the lemmas of Secs.~\ref{sec:direct}--\ref{sec:removing} on fresh test states. Every check reports the fraction of its tolerance that it uses; none uses more than $20\%$.
